\documentclass[11pt,a4paper]{article}
\usepackage[utf8]{inputenc}
\usepackage[T1]{fontenc}
\usepackage[spanish]{babel}
\usepackage{hyperref}
\usepackage{geometry}
\geometry{margin=2.5cm}
\usepackage{parskip}

\title{328 aventarse sin condiciones daña}
\author{Alejandro Valenzuela Rivera\\\small Universidad Autónoma de Nuevo León}
\date{\today}
\begin{document}
\maketitle

\begin{quote}
\textbf{Sesgo declarado.} Este trabajo se sostiene desde el sesgo declarado de su autor, y no finge neutralidad. Lo que afirma se puede tumbar, y ese es el punto: el falsificador está escrito, no escondido.
\end{quote}

\section*{Resumen}
El contacto reduce el prejuicio, y solo con cuatro condiciones: estatus igual, metas comunes, cooperación y respaldo de la autoridad. Sin ellas, no es neutro: daña.

\section{Por qué se sostiene}
La regla se sostiene mientras alguien la pueda contestar con lo que Allport, 713 muestras de 515 estudios, más de 250 000 sujetos muestra.

\section{De dónde viene}
Tesis/publicable/casos.md.

\section{El caso}
Allport: 713 muestras de 515 estudios, y el encuentro sin condiciones aumenta el prejuicio.

\section{Método}
El contacto negativo empeora más de lo que el positivo mejora: bad is stronger than good.

\section{Límites, y qué lo tumba}
Un caso donde juntar sin condiciones reduzca el prejuicio.

\section{Lo que queda abierto}
\textquestiondown{}El encuentro tuvo las cuatro condiciones, o solo se juntaron?

\section*{Cómo se leyó esta tesis}
Esto se escribió en el repositorio man-made-empathy, del que este candado es una pieza: candados/06 del caso/328 aventarse sin condiciones daña.md. La máquina que lo compuso lo hizo desde este mismo archivo, y no agregó nada que no estuviera escrito. Si algo de esto es falso, se dice en el mismo sitio, y se corrige con fecha.

\end{document}
